\documentclass[11pt,a4paper]{article}
\usepackage[margin=25mm]{geometry}
\usepackage{amsmath,amssymb,amsthm,booktabs}
\newtheorem{lemma}{Lemma}\newtheorem{theorem}{Theorem}
\title{Four order 14 counts from projective overgroups}
\author{Human Compression project}\date{}
\begin{document}\maketitle
Let $G=O'N$, and let $x_J(p,q)$ count ordered pairs $(r,s)\in J^2$
with $|r|=p$, $|s|=q$ and $|rs|=2$. Write $e_J(p,q)$ for the
number of these pairs generating $J$. We treat
$p\in\{3,6,8,12,16\}$ and $q=14$. The last case is an independent
check of an already recovered entry. All subgroup row numbers refer
to the complete acquired roster. Its completeness and the acquired
character tables remain explicit borrowed premises.

\begin{lemma}\label{orthogonal}
For $M=\operatorname{PSL}_3(7):\langle\tau\rangle$, with $\tau$
inverse transpose, the centralizer of an outer involution is
$\operatorname{PGL}_2(7)\times C_2$, of order 672.
\end{lemma}
\begin{proof}
We reuse the accepted symmetric-form argument for this same
projective model. If $t=[A]\tau$ is an involution, then
$A=\zeta A^T$ with $\zeta^3=1$. Transposing gives $\zeta^2=1$,
so $A$ is symmetric. The inner centralizer is the projective
similitude group of a nondegenerate ternary form. A similitude's
multiplier is a square, because $(\det h)^2=\mu^3$; each projective
class therefore has a unique isometric representative of determinant
one. These representatives lie in $\operatorname{SL}_3(7)$.
Consequently the inner centralizer is $\operatorname{SO}_3(7)$,
and the full centralizer is its direct product with $\langle t\rangle$.

For completeness, we identify this orthogonal group explicitly.
A nondegenerate ternary form over $\mathbb F_7$ is isotropic.
After diagonalizing it, the sets of values $aX^2$ and $-c-bY^2$
each have four elements, so intersect. A Witt basis now makes
the form similar to $Q(a,b,c)=a^2+bc$. Conjugation of traceless
two by two matrices preserves $-\det=Q$ and is faithful modulo
scalar matrices. Its determinant on this three-dimensional space
is one: this holds directly for diagonal matrices and elementary
unipotents, which generate $\operatorname{GL}_2(7)$. Thus it embeds
$\operatorname{PGL}_2(7)$, of order 336, in $\operatorname{SO}(Q)$.

There are exactly eight isotropic lines, represented by
$[a,1,-a^2]$ for $a\in\mathbb F_7$, and $[0,0,1]$.
Any three distinct such lines form a basis, and any two have
nonzero bilinear product. These statements follow immediately
from the displayed coordinates (the determinant for three finite
parameters is a nonzero Vandermonde determinant).
An isometry fixing three such lines has eigenvalues $\lambda_i$
satisfying $\lambda_i\lambda_j=1$ for $i\ne j$. They are all the
same sign; determinant one forces the positive sign. The action
on ordered triples of distinct isotropic lines is therefore free,
so $|\operatorname{SO}(Q)|\leq8\cdot7\cdot6=336$.
The embedding is an isomorphism, proving the assertion.
\end{proof}

The following small counts have complete literal certificates.
Let $B$ be the order 8232 flag normalizer from the accepted order 28
argument, let $H$ be the order 1344 centralizer of an inner involution
in $M$, and let $K=\operatorname{PGL}_2(7)\times C_2$.
\[
\begin{array}{c|rrrrr}
p&3&6&8&12&16\\\hline
x_B(p,14)&0&0&0&0&0\\
x_H(p,14)&0&1344&672&672&2688\\
x_K(p,14)&336&1008&1344&0&0.
\end{array}                                                    \tag{1}
\]
The reader rechecks the entire retained sets of 8232, 1344 and 168
permutations on 114 letters, their generator closures and their
order 14 conjugacy classes. There are three such classes in $B$,
of sizes 294, 588 and 588, and one in $H$, of size 48. It checks
every possible partner for each representative and weights by the
class size. It also verifies an order 28 square root in $H$;
conjugacy gives a root for every order 14 element of $H$.
For $K$, the reader independently constructs the complete projective
action of all 336 nonsingular two by two matrices modulo scalars
on eight projective points, and enumerates its direct product with
$C_2$. The first factor has no element of order 14; all order 14
elements of $K$ therefore have nontrivial second coordinate.

\begin{lemma}\label{inner}
Each element $s\in G$ of order 14 lies in a unique row 7 maximal
subgroup $C=C_G(s^7)$. There is exactly one projective maximal
subgroup of each of rows 2 and 3 in which $s$ is inner. Call them
$M_2,M_3$. Their intersection is $B$, and
$M_i\cap C$ is an order 1344 subgroup $H$ in row 65 or 66,
respectively.
\end{lemma}
\begin{proof}
The unique cyclic order 14 class, row 534, occurs in cyclic order
28, row 466. The primary ATLAS table gives $|C_G(s)|=28$ for the
sole class 14A. Thus $C_G(s)$ is cyclic of order 28 and has the
unique involution $t=s^7$. Every row 7 maximal is an involution
centralizer, so precisely one can contain $s$.

Suppose $s$ is inner in a projective maximal $M$. Then $t$ is
inner too, and the accepted block calculation gives
$C_M(t)\cong H$. The square-root assertion following (1) shows
that $M$ contains the full cyclic centralizer $C_G(s)$.
Conversely, in any projective maximal containing this cyclic
order 28 subgroup, its subgroup of squares is inner, because
$M/\operatorname{PSL}_3(7)\cong C_2$. The accepted order 28
overgroup theorem gives exactly one such $M$ in each projective
class, their intersection $B$, and their intersections with $C$
as the indicated block centralizers. This proves all the assertions.
\end{proof}

\begin{lemma}\label{outer}
In a fixed projective maximal subgroup $M$, every nongenerating
pair of our five types whose order 14 element $s$ is outer lies
in the single subgroup $C_M(s^7)\cong K$. The total number of
such pairs in $M$ is $5586x_K(p,14)$.
\end{lemma}
\begin{proof}
The complete maximal list of $M$, including actual subgroup counts,
is as follows. The two rows shown for a type correspond to the
two global projective classes.
\[
\begin{array}{c|rrrrrrr}
\text{rows}&4/5&22&65/66&110/111&250/252&251/253&285\\
\text{order}&1876896&8232&1344&672&144&144&114\\
\text{number}&1&456&2793&5586&26068&26068&32928.
\end{array}                                                     \tag{2}
\]
The socle cannot contain outer $s$, and the order 8232 subgroups
have no relevant pair by (1). An order 1344 inner-involution
centralizer cannot contain $s$: its central inner involution
would commute with $s$ and hence equal the outer involution $s^7$.
Orders 144 and 114 exclude an element of order seven.

The centralizer $C_M(s^7)$ from Lemma~\ref{orthogonal} is itself
maximal. Its order 672 excludes all candidates in (2) except the
socle and the order 1344 subgroup; the former excludes its outer
involution, and the latter is excluded by any outer order 14
element of $K$, as just argued. It is consequently a representative
of the unique maximal class of order 672 in (2). Every subgroup
in that class is an outer-involution centralizer. In any such
subgroup containing $s$, the central involution must be $s^7$.
Thus a pair occurs in only one of these subgroups within $M$.
Conversely every pair counted in $K$ is proper in $M$ and has
outer $s$. The count 5586 in (2) now proves the formula.
\end{proof}

\begin{theorem}
For $p\in\{3,6,8,12,16\}$,
\[
 e_G(p,14)=x_G(p,14)-\sum_{M}[G:M]x_M(p,14)
       +2\cdot342868680\bigl(x_H(p,14)+2x_K(p,14)\bigr),       \tag{3}
\]
where the sum has one term for each of the 13 maximal subgroup
classes. The resulting ordered automorphism-orbit counts are
\[
\begin{array}{c|rrrrr}
p&3&6&8&12&16\\\hline
e_G(p,14)/|\operatorname{Aut}(G)|&10&597&3056&1391&12504.
\end{array}
\]
\end{theorem}
\begin{proof}
The acquired character tables give nonzero maximal contributions
only in rows 2, 3 and 7. Start with their full sum. In each
projective maximal, remove the proper pairs with outer $s$ as in
Lemma~\ref{outer}; these pairs are already covered by $C_G(s^7)$.
There are 122760 projective maximals in each class, so this removal
is
\[
2\cdot122760\cdot5586x_K=4\cdot342868680x_K.
\]
This is an incidence count over actual projective maximals, not
an assertion that an order 672 subgroup has a unique maximal
overgroup. In fact each row 110 or 111 subgroup belongs to two
maximals of its projective class.

The remaining projective pairs are either inner or generate the
whole projective maximal. A generating pair cannot lie in a second
maximal subgroup of $G$. For an inner pair, Lemma~\ref{inner}
gives its canonical projective maximals and its unique row 7
maximal. The two projective contributions have no common pair,
since $x_B=0$ in (1). Their overlaps with row 7 are exactly the
order 1344 groups in rows 65 and 66. Each has class length
342868680; each contains only inner order 14 elements, since its
central inner involution must be $s^7$. The accepted unique
overgroup theorem for these intersections therefore gives the
remaining correction $2\cdot342868680x_H$. No triple overlap
survives, again because $x_B=0$. This proves (3).

Substitution of (1) and the acquired character evaluations gives
the stated integers after division by $|\operatorname{Aut}(G)|=2|G|$.
This action on generating pairs is free. There is no division by
duality. All five answers agree with the fixed source; the fifth
was already known from the order 16 calculation.
\end{proof}

\paragraph{Attribution and limits.}
The numerical assertions are from Connor and Leemans,
\emph{Algorithmic enumeration of regular maps}, Table~1. We borrow
the complete subgroup roster, the primary ATLAS centralizer data
and the acquired character evaluations, with their existing
provenance. The block centralizer and order 28 intersection theorems
are accepted earlier components. The present argument separates
inner and outer order 14 contributions and proves the required
orthogonal identification explicitly. The literal certificates
involve only the displayed small groups. We make no priority claim
for the elementary group theory, and give no additional chirality
or whole-original completion credit here.
\end{document}
